\subsubsection{The Strong Maximum Principle}\label{subsec:strong-maximum-principle}

Theorem~\ref{thm:weak-maximum-principle} states that $u$ assumes its maximum at the boundary. However, $u$ may assume its maximum at many points, and therefore the theorem does not rule out the possibility that some of these points are interior. The strong maximum principle states that this is impossible, unless $u$ is a constant. For the proof, we shall need the following lemma, which is interesting in its own right.

\begin{figure}[H]
\centering
\includegraphics[width=.78\linewidth]{figures/boundary-point-interior-ball-source.png}
\caption{}\label{fig:boundary-point-interior-ball}
\end{figure}

\begin{Lemma}
\label{lem:elliptic-boundary-point}
Suppose that $\Omega$ lies on one side of $\partial\Omega$. Assume $Lu\geq0$, and let $\bm{x}_0$ be a point on $\partial\Omega$ such that $u(\bm{x}_0)>u(\bm{x})$ for every $\bm{x}\in\Omega$. Also assume that, in a neighborhood of $\bm{x}_0$, $\partial\Omega$ is a $C^2$-surface and that $u$ is differentiable at $\bm{x}_0$. Moreover, suppose that either
\begin{enumerate}
\item $c=0$,
\item $c\leq0$ and $u(\bm{x}_0)\geq0$, or
\item $u(\bm{x}_0)=0$
\end{enumerate}
Then $\partial u/\partial n(\bm{x}_0)>0$, where $\partial u/\partial n$ denotes the derivative in the direction of the outer normal to $\partial\Omega$.
\end{Lemma}

\begin{Proof}
Since $\partial\Omega$ was assumed $C^2$, we can choose (see Problem~\ref{ex:interior-tangent-ball}) a ball $B_R(\bm{y})$ such that $B_R(\bm{y})\subset\Omega$ and $\bm{x}_0\in\partial B_R(\bm{y})$. Here $R$ and $\bm{y}$ denote the radius and center of the ball.

For $0\leq r=|\bm{x}-\bm{y}|\leq R$, define
\begin{equation}\label{eq:boundary-point-exponential-barrier}
v(\bm{x})=\exp(-\alpha r^2)-\exp(-\alpha R^2).
\end{equation}
We find
\begin{equation}\label{eq:boundary-point-barrier-operator}
\begin{aligned}
Lv(\bm{x})&=\exp(-\alpha r^2)\bigl[4\alpha^2a_{ij}(x_i-y_i)(x_j-y_j)\\
&\qquad-2\alpha(a_{ii}+b_i(x_i-y_i))\bigr]+cv.
\end{aligned}
\end{equation}
Now let $A=B_R(\bm{y})\cap B_{R'}(\bm{x}_0)$, with $R'$ chosen small. For large enough $\alpha$, we have $Lv>0$ in $A$. Moreover, if we choose $\epsilon>0$ small enough then $u-u(\bm{x}_0)+\epsilon v\leq0$ on $\partial A\cap\partial B_{R'}(\bm{x}_0)$, and also on $\partial A\cap\partial B_R(\bm{y})$, where $v=0$. Thus we find $L(u-u(\bm{x}_0)+\epsilon v)\geq-cu(\bm{x}_0)\geq0$ in $A$ and $u-u(\bm{x}_0)+\epsilon v\leq0$ on $\partial A$.

If $c\leq0$, the weak maximum principle (Corollary~\ref{cor:weak-maximum-nonpositive-zero-order}) implies that $u-u(\bm{x}_0)+\epsilon v\leq0$ throughout $A$. We take the normal derivative at $\bm{x}_0$ and obtain
\begin{equation}\label{eq:boundary-point-positive-normal-derivative}
\dfrac{\partial u}{\partial n}(\bm{x}_0)\geq-\epsilon\dfrac{\partial v}{\partial n}(\bm{x}_0)=2\alpha\epsilon R\exp(-\alpha R^2)>0,
\end{equation}
which implies the lemma.

If $u(\bm{x}_0)=0$, then, by assumption, $u$ is negative in $\Omega$. Now let $c^+(\bm{x})=\max(0,c(\bm{x}))$. We find that $(L-c^+)u=Lu-c^+u\geq Lu\geq0$, and hence we can apply the argument above with $L-c^+$ in place of $L$.
\end{Proof}

\begin{Remark}
\label{rem:boundary-one-side-jordan}
Since $\Omega$ is assumed to be connected, it can be shown that $\Omega$ lies on one side of $\partial\Omega$ if $\Omega$ is bounded and $\partial\Omega$ is globally smooth. (This is a multidimensional generalization of the Jordan curve theorem.) For a proof of this see, e.g.,~\cite{Mas}.
\end{Remark}

\begin{Remark}
\label{rem:boundary-point-semidefinite}
Lemma~\ref{lem:elliptic-boundary-point} still holds if the matrix $a_{ij}$ is only positive semidefinite and $\bm{n}$ is not in the nullspace.
\end{Remark}

As a consequence of Lemma~\ref{lem:elliptic-boundary-point}, we obtain the following strong maximum principle.

\begin{Theorem}
\label{thm:strong-maximum-principle}
Assume $Lu\geq0$ ($Lu\leq0$) in $\Omega$ (not necessarily bounded) and assume that $u$ is not constant. If $c=0$, then $u$ does not achieve its maximum (minimum) in the interior of $\Omega$. If $c\leq0$, $u$ cannot achieve a non-negative maximum (non-positive minimum) in the interior. Regardless of the sign of $c$, $u$ cannot be zero at an interior maximum (minimum).
\end{Theorem}

\begin{Proof}
Assume that $u$ assumes its maximum $M$ at an interior point and let $\Omega^-=\Omega\cap\{u<M\}$. If $\Omega^-$ is not empty, then $\partial\Omega^-\cap\Omega$ is not empty. Let $\bm{y}$ be a point in $\Omega^-$ that is closer to $\partial\Omega^-$ than to $\partial\Omega$ and let $B$ be the largest ball contained in $\Omega^-$ and centered at $\bm{y}$. Let $\bm{x}_0$ be a point on $\partial B\cap\partial\Omega^-$. Then we can apply the previous lemma to $B$. We conclude that $\nabla u$ is nonzero at $\bm{x}_0$, contradicting the assumption that $u$ assumes its maximum there.
\end{Proof}
